% page 109 of the facsimile (image p-108 of the scan)
% block 182.9 x 250.6 mm ; left margin 19.1 mm
\pgc{-2.540mm}{0.000mm}{
\l{}
\l{}
\l{\cel{59}lOl}
\l{}
\l{}
\l{\cel{5}\sou{delfino}. I. (\sou{historio di Francia}).Titulo di la filiulo seniora di}
\l{\cel{8}la reji di Francia. - II. (\sou{zool}.) Cetaceo suflera, mikra, di}
\l{\cel{8}qua la kranio esas konvexa, quan la fabli antiqua ed olima re-}
\l{\cel{8}prezentas kom amika a la homo. - L.\cel{1}\sou{delphinus}.\cel{2}- DEFIRS.}
\l{}
\l{\cel{5}\sou{delfinelo}.\cel{2}(\sou{bot}.) Nomo di ranunkulaceo, planto astriktiva.}
\l{\cel{8}- L.\cel{1}\sou{delphinium}.- F.}
\l{}
\l{\cel{5}\sou{deliberar}.\cel{2}(\sou{netrans., pri}). Ponderar en su ipsa, o kun altra}
\l{\cel{8}personi, pri problemo solvenda o rezolvo agenda, la motivin}
\l{\cel{8}\textquotedbl{}por\textquotedbl{} e la motivi \textquotedbl{}kontre\textquotedbl{} la decido o rezolvo. - DEFIS.}
\l{}
\l{\cel{5}\sou{delico.}\cel{2}Juo delikatega, bonega, ecelanta. - EFIS.}
\l{}
\l{\cel{5}\sou{delikata}. Qua esas extreme dina, tenua, pura : di qua la gusto-}
\l{\cel{8}sentiveso igas nefacila la satisfaco. - DEFIRS.}
\l{}
\l{\cel{5}\sou{deliktar}. (\sou{netrans.})\cel{2}Agar tale ke la lego esas violacita. -}
\l{\cel{8}(En\cel{1}\sou{Francia}) Opoze a krimino ed a kontravenco : violaco di la}
\l{\cel{8}lego qua efektigas puniseso segun la delikto-kodexo. - DEFIRS.}
\l{}
\l{\cel{5}\sou{deliquecar}. (\sou{netrans.})\cel{2}(\sou{Kemio)}.\cel{2}Qua aquiris la proprajo absorbar}
\l{\cel{8}la vaporo di aero humida e, pro lo, dissolvesar. - DEFIS.}
\l{}
\l{\cel{5}\sou{delirar}.\cel{2}(\sou{netrans.})\cel{2}I. Standar tale ke, pro febro, ebrieso, e c.,}
\l{\cel{8}la idei esas acidente nekoheranta ed ecitegita. - II. Havar sua}
\l{\cel{8}spirito en ta stando di exalteso plu o min violentoza qua}
\l{\cel{8}karakterizas alienaco. - III. Havar sua spirito en tala stando}
\l{\cel{8}di exalteso violentoza, ke lu esas obediar raciono. - DEFIRS.}
\l{}
\l{\cel{5}\sou{deltoido}. (\sou{Anat}.) Muskulo triangulatra qua, per juntar la humero}
\l{\cel{8}a la klavikulo ed a la skapulo, servas por elevar ed abasar}
\l{\cel{8}la brakio.}
\l{}
\l{\cel{5}\sou{demagogo}. La homo qua esforcas igar su povoza a turbo, per flatar}
\l{\cel{8}lu. - DEFIS.}
\l{}
\l{\cel{5}\sou{demagogio}. I. Sistemo di politiko segun qua onu flatas turbo. -}
\l{\cel{8}II. Stato ube la povo statala abandonesas a turbo. - DEFIRS.}
\l{}
\l{\cel{5}\sou{demandar}. (\sou{trans., ulo, ad ulu}).- I. Savigar ke onu volas havar}
\l{\cel{8}(ulo). - II. Savigar da ulu ke onu expektas (ulo) de lu. - FIS.}
\l{}
\l{\cel{5}demarshar. (\sou{netrans.})\cel{2}Ago sociala per qua persono volas influar}
\l{\cel{8}altra por obtenar ulo. Demarsho implikas sempre la ideo di}
\l{\cel{8}\textquotedbl{}pazi\textquotedbl{} en la socio por atingar ula skopo, por obtenar irga}
\l{\cel{8}helpo o favoro. - F.}
\l{}
\l{\cel{5}dementa. (\sou{Patol}.) Malada de grava desordineso di sua raciono.-}
\l{\cel{8}DEFIS.}
\l{}
\l{\cel{5}dementiar. (\sou{trans.})\cel{2}I. Kontredicar (ulu) kom nedicinta lo exakta.}
\l{\cel{8}II. Kontredicar (to quon dicas ulu) kom neexakta. - FIS.}
}
